%========================
% Single-column CV (Hope Donglo) -- University of Ghana, Department of Physics
%========================
\documentclass[10pt,a4paper]{article}

\usepackage[scaled]{helvet}
\renewcommand{\familydefault}{\sfdefault}

\usepackage[utf8]{inputenc}
\usepackage[T1]{fontenc}
\usepackage[english]{babel}
\usepackage[a4paper, top=1.5cm, bottom=1.5cm, left=1.5cm, right=1.5cm]{geometry}
\usepackage{enumitem}
\usepackage{titlesec}
\usepackage{ragged2e}
\usepackage{xcolor}
\usepackage{hyperref}
\usepackage{xparse}
\usepackage{amsmath}

\definecolor{maincolor}{rgb}{0.5,0.05,0}

\newcommand{\cvsection}[1]{%
  \par\vspace{8pt}%
  \noindent{\large\bfseries\textcolor{maincolor}{#1}}\par
  {\color{maincolor}\rule{\linewidth}{0.6pt}}%
  \vspace{4pt}%
}

\newcommand{\cvsubsection}[1]{%
  \par\vspace{4pt}%
  \noindent{\normalsize\bfseries\textcolor{maincolor}{#1}}%
  \vspace{2pt}%
}

\newlength{\tgutter}\setlength{\tgutter}{1.5em}
\newlength{\tspace}\setlength{\tspace}{1pt}
\newsavebox{\tbox}

\NewDocumentCommand{\timelineitem}{s m m m}{%
  \par\noindent
  \sbox{\tbox}{%
    \begin{minipage}[t]{\dimexpr\linewidth-\tgutter\relax}
      {\small\textbf{#2}}\\[-1pt]
      {\color{maincolor}\textit{\small #3}}\\[-1pt]
      {\footnotesize
        \begin{itemize}[leftmargin=1.2em,itemsep=1pt,topsep=3pt,
                        label=\textbullet]
          #4
        \end{itemize}
      }
    \end{minipage}%
  }%
  \makebox[0pt][r]{%
    \begin{minipage}[t]{\tgutter}
      \centering
      {\small\textbullet}\\[-0.3ex]
      \IfBooleanTF{#1}{}{\rule{0.4pt}{\dimexpr\ht\tbox+\dp\tbox+
        \tspace-0.8ex\relax}}%
    \end{minipage}%
  }%
  \usebox{\tbox}\vspace{\tspace}%
}

\newlength{\dgutter}\setlength{\dgutter}{1.5em}
\newlength{\dspace}\setlength{\dspace}{2pt}
\newsavebox{\dbox}

\NewDocumentCommand{\diploma}{s m m m}{%
  \par\noindent
  \sbox{\dbox}{%
    \begin{minipage}[t]{\dimexpr\linewidth-\dgutter\relax}
      {\small\textbf{#2}}\\[-1pt]
      {\color{maincolor}\textit{\small #3}}\\[-2pt]
      {\footnotesize
        \begin{itemize}[leftmargin=1.2em,itemsep=2pt,topsep=2pt,
                        label=\textbullet]
          #4
        \end{itemize}
      }
    \end{minipage}%
  }%
  \makebox[0pt][r]{%
    \begin{minipage}[t]{\dgutter}
      \centering
      {\small\textbullet}\\[-0.3ex]
      \IfBooleanTF{#1}{}{\rule{0.4pt}{\dimexpr\ht\dbox+\dp\dbox+
        \dspace-0.8ex\relax}}%
    \end{minipage}%
  }%
  \usebox{\dbox}\vspace{\dspace}%
}

\setlist[itemize]{noitemsep, topsep=2pt, leftmargin=1.4em}
\pagenumbering{arabic}
\setlength{\parskip}{1pt}
\setlength{\parindent}{0pt}
\hypersetup{colorlinks=true, linkcolor=black,
            urlcolor=black, citecolor=black,
            pdftitle={Hope Donglo, PhD: Curriculum Vitae},
            pdfauthor={Hope Donglo}}

\begin{document}
\selectlanguage{english}

\noindent\colorbox{maincolor}{%
  \begin{minipage}{\dimexpr\linewidth-2\fboxsep\relax}
    \vspace{8pt}
    \begin{center}
      {\Huge\textbf{\textcolor{white}{Hope Donglo, PhD}}\par}
      \vspace{4pt}
      {\normalsize\textcolor{white}{%
        Lecturer, Department of Physics \enspace\textbar\enspace
        University of Ghana, Legon}\par}
      \vspace{5pt}
      \hypersetup{urlcolor=white}%
      {\footnotesize\textcolor{white}{%
        Legon, Accra, Ghana \enspace\textbar\enspace
        020\,48\,10\,352 \enspace\textbar\enspace
        \href{mailto:hdonglo@ug.edu.gh}{hdonglo@ug.edu.gh}
        \enspace\textbar\enspace
        \href{mailto:hkdonglo@gmail.com}{hkdonglo@gmail.com}
        \enspace\textbar\enspace
        \href{https://linkedin.com/in/hopedonglo}{%
          linkedin.com/in/hopedonglo}%
      }}
    \end{center}
    \vspace{8pt}
  \end{minipage}%
}
\hypersetup{urlcolor=black}%
\vspace{6pt}

\cvsection{Professional Profile}

{\footnotesize
Theoretical nuclear physicist (PhD, Universit\'{e} de Caen Normandy;
Marie Sk\l{}odowska-Curie -- COFUND Fellow) specialising in nuclear
reaction dynamics and superheavy element synthesis, with research
experience at CEA/GANIL, INFN Legnaro, and the Universidad Aut\'{o}noma
de Madrid. Develops theoretical models of heavy-ion reactions and
superheavy element synthesis, employing stochastic methods, uncertainty
quantification, and scientific computing to investigate reaction
mechanisms and improve predictive nuclear theory. Contributes to
development and validation of the KEWPIE3 code.
}

\vspace{4pt}
{\footnotesize
\begin{tabular}{@{}p{4.6cm}p{10.5cm}@{}}
\textbf{Rank:} & Lecturer \\
\textbf{Highest Qualification:} & PhD, Theoretical Nuclear Physics (Universit\'{e} de Caen Normandy, 2024) \\
\textbf{Date of Appointment (UG):} & 1 August 2026 \\
\textbf{Date of Assumption of Duty:} & 17 August 2026 \\
\textbf{Area of Specialisation:} & Nuclear reaction theory; superheavy element synthesis; uncertainty quantification \\
\end{tabular}
}

\cvsection{Education}

\diploma
{10/2021 -- 12/2024 \quad Ph.D.\ in Theoretical Nuclear Physics}
{University of Caen Normandy, France}
{
  \item \textit{Dissertation:} Study of Reaction Mechanisms for the
        Synthesis of Superheavy Elements
  \item Marie Sk\l{}odowska-Curie -- COFUND Fellow; research conducted at
        CEA/GANIL (France)
  \item \textit{Supervisor:} Dr.\ Dieter Ackermann (GANIL/CEA)
}

\diploma
{09/2018 -- 10/2020 \quad Joint Master's in Nuclear Physics
(Erasmus Mundus)}
{Universities of Seville, Autonomous Madrid, Padua, and Caen}
{
  \item Competitive international scholarship; training across four
        European universities
  \item \textit{Thesis:} Nuclear Structure Studies with
        Symmetry-Conserving Configuration Mixing Method
  \item \textit{Supervisor:} Dr.\ Tom\'{a}s Ra\'{u}l Rodr\'{i}guez
        Frutos
}

\diploma
{09/2012 -- 12/2016 \quad Bachelor of Science in Physics}
{University of Ghana, Legon --- thesis conducted at GAEC/NNRI}
{
  \item \textit{Thesis:} Evaluation of Safety Status of the Ghana
        Research Reactor-1 (GHARR-1)
  \item \textit{Supervisors:} Dr.\ Allison F.\ Hughes (UG Physics)
        and Dr.\ Joseph K.\ Gbadago (GAEC/NNRI)
}

\cvsection{Academic and Research Appointments}

\timelineitem
{08/2026 -- \quad Lecturer, Department of Physics}
{College of Basic and Applied Sciences (CBAS), University of Ghana, Legon}
{
  \item Course allocation to be determined
}

\timelineitem
{10/2021 -- 12/2024 \quad Doctoral Researcher,
Theoretical Nuclear Physics}
{CEA/GANIL, Caen, France}
{
  \item Modelled nuclear reaction mechanisms using stochastic
        simulation methods (Langevin dynamics), covering the full
        synthesis pathway for superheavy element production
  \item Built and validated HPC simulation pipelines (KEWPIE2/KEWPIE3
        framework); predicted production cross-sections for
        $Z=119$ and $Z=120$ nuclei
  \item Applied Bayesian optimisation and machine learning methods
        for nuclear model parameter estimation and uncertainty
        quantification
  \item Performed systematic benchmarking of competing theoretical
        models with rigorous uncertainty analysis
}

\timelineitem
{03/2020 -- 09/2020 \quad Graduate Research Assistant}
{Autonomous University of Madrid (UAM), Spain}
{
  \item Simulated 27 argon isotopes ($A=32$--58) using the
        Symmetry-Conserving Configuration Mixing method; mapped
        structural evolution and shape coexistence across the
        isotopic chain
}

\timelineitem
{09/2019 -- 01/2020 \quad Graduate Research Assistant}
{GANIL, Caen, France}
{
  \item Investigated nuclear deformation dynamics using the
        Geometric Collective Model
  \item Identified $\gamma$-unstable shape coexistence and performed
        numerical simulations of quadrupole dynamics
}

\timelineitem*
{03/2019 -- 06/2019 \quad Graduate Research Assistant}
{National Institute for Nuclear Physics (INFN), Legnaro, Italy}
{
  \item Optimised radiation detectors (EJ-301, LaBr$_3$, BaF$_2$)
        for neutron/gamma discrimination; achieved 0.01\,ns/channel
        TOF calibration and ${<}5\%$ LaBr$_3$ energy resolution
        deviation
  \item Developed automated data acquisition and analysis pipelines
        for neutron time-of-flight spectroscopy
}

\cvsection{Publications}

{\footnotesize
\begin{itemize}[leftmargin=1.4em, itemsep=6pt, topsep=2pt,
                label=\textbullet]

  \item \textbf{Peer-Reviewed}\\
  D.\ Ackermann, F.\ Agert, A.\ Bahini, D.\ Boilley, \textbf{H.\ Donglo},
  S.\ Kumar, J.\ Piot, H.\ Savajols, Ch.\ Stodel, and the S$^3$ Collaboration (2026).
  ``Probing SHN Structure and Stability -- synthesis and decay studies at S$^3$.''
  \textit{Acta Physica Polonica B}, Vol.\ 19 (2026).

  \item \textbf{Preprint (HAL)}\\
  \textbf{H.\ Donglo} and T.~R.\ Rodr\'iguez Frutos.
  ``Nuclear Structure of Even-Even Argon Isotopes within the Axial
  Symmetry Conserving Configuration Mixing Framework.''
  Deposited on HAL (hal-05677054).

  \item \textbf{PhD Thesis}\\
  \textbf{H.\ Donglo} (2024).
  \textit{Study of Reaction Mechanisms for the Synthesis of
  Superheavy Elements.}
  University of Caen Normandy.

\end{itemize}
}

\cvsection{Ongoing Research Projects}

{\footnotesize
\begin{itemize}[leftmargin=1.4em, itemsep=6pt, topsep=2pt,
                label=\textbullet]

  \item \textbf{KEWPIE3 Development and Validation}\\
  Validation and publication of KEWPIE3, the new version of the KEWPIE2 code
  for fusion-evaporation reaction cross-section predictions in superheavy
  element synthesis. Contributes to improved compound nucleus formation
  probability calculations using Langevin-based stochastic approaches.

  \item \textbf{In Preparation}\\
  \textbf{H.\ Donglo}, F.\ Agert, D.\ Boilley, T.\ Cap, M.\ Kowal, and D.\ Ackermann.
  ``Revisiting Fusion Dynamics in the Fusion-by-Diffusion Model for
  Fusion--Evaporation Production Cross Sections.''
  Targeting \textit{Journal of Physics G}.

  \item \textbf{In Preparation}\\
  \textbf{H.\ Donglo}, D.\ Boilley, F.\ Cap, and D.\ Ackermann.
  ``Constraining the Compound Nucleus Formation Probability in
  Superheavy Element Synthesis through Systematic Uncertainty
  Analysis.'' Targeting \textit{Physical Review C}.

\end{itemize}
}

\cvsection{Research Interests}

{\footnotesize
Nuclear reaction theory (fusion-evaporation, superheavy element
synthesis); nuclear structure (collective and mean-field models);
Langevin dynamics and stochastic simulation; Bayesian inference and
uncertainty quantification for nuclear model parameters; machine
learning methods for nuclear data analysis; simulation-assisted
nuclear safety analysis; nuclear fission.
}

\cvsection{Teaching and Supervision}

\cvsubsection{Teaching Experience}

{\footnotesize
\textbf{2022--2024 \quad BSc Intern Co-Supervisor}\\
GANIL/CEA, Caen, France\\
Supervised final-year interns on nuclear structure projects; guided
Python-based simulation workflows, uncertainty analysis, and scientific writing.

\vspace{3pt}
\textbf{09/2016--08/2017 \quad Undergraduate Teaching Assistant}\\
University of Ghana, Department of Physics, Legon\\
Supervised laboratory sessions for 150+ undergraduates across
experimental physics modules; designed structured data analysis
exercises and provided individual mentoring in scientific computing.
Supported curriculum delivery alongside academic staff.

\vspace{3pt}
\textbf{2017--2021 \quad Physics and Mathematics Tutor (Intermittent, Part-time)}\\
Greater Accra, Ghana\\
Delivered secondary-level instruction in physics and mathematics;
adapted teaching approaches to diverse learner backgrounds and
ability levels.
}

\cvsubsection{Teaching Areas}

{\footnotesize
\begin{itemize}[leftmargin=1.4em, itemsep=2pt, topsep=2pt,
                label=\textbullet]
  \item Introductory Physics and Laboratory Instruction
  \item Nuclear and Particle Physics
  \item Quantum Mechanics and Quantum Field Theory
  \item Mechanics, Thermodynamics, and Electromagnetism
  \item Computational and Scientific Programming Methods for Physics
\end{itemize}
}

\cvsubsection{Teaching Philosophy}

{\footnotesize
Committed to inquiry-based, student-centred physics instruction that
integrates conceptual reasoning, laboratory practice, and
computational problem-solving. Championing inclusive, research-informed STEM education.
}

\cvsection{Fellowships and Awards}

{\footnotesize
\begin{itemize}[leftmargin=1.4em, itemsep=3pt, topsep=2pt,
                label=\textbullet]
  \item \textbf{Marie Sk\l{}odowska-Curie Actions -- COFUND Fellow}
        (2021--2024) --- European Commission, fully funded doctoral
        fellowship
  \item \textbf{Erasmus Mundus Scholar} (2018--2020) ---
        European Commission, joint MSc scholarship
\end{itemize}
}

\cvsection{Professional Development and Training}

{\footnotesize
\textbf{10/2025 -- 09/2026 \quad Data and ML Engineering
Certification} (Bac+5, RNCP Level 7) ---
\textit{OpenClassrooms, France Travail-funded}\\
Professional certification in scalable data pipelines, MLOps, and
machine learning applications, with relevance to computational
physics instruction and research methods teaching.

\vspace{4pt}
\textbf{2024 \quad Research Integrity Certification} ---
Universit\'{e} de Bordeaux
}

\cvsection{Conferences and Scientific Presentations}

{\footnotesize
\textbf{Seminars and Conference Presentations}
\begin{itemize}[leftmargin=1.4em, itemsep=2pt, topsep=2pt,
                label=\textbullet]
  \item Study of Reaction Mechanisms for the Synthesis of Superheavy
        Elements. \textit{29th Nuclear Physics Conference}, Poland,
        September 2023.
\end{itemize}

\vspace{2pt}
\textbf{Workshop Participation}
\begin{itemize}[leftmargin=1.4em, itemsep=1pt, topsep=1pt,
                label=\textbullet]
  \item GRIT Workshop --- Silicon Detectors for Low-Energy Nuclear
        Physics, GANIL, Caen, France, 2023.
  \item Workshop on Targets and Ion Sources, GANIL,
        Caen, France, 2023.
  \item LISE Workshop --- Nuclear Structure and Reactions, GANIL,
        Caen, France, 2022.
  \item Physics with SPIRAL2 Heavy Ion Beams, GANIL, Caen, France,
        October 2022.
\end{itemize}

\vspace{2pt}
\textbf{Schools and Training}
\begin{itemize}[leftmargin=1.4em, itemsep=1pt, topsep=1pt,
                label=\textbullet]
  \item \textbf{CEA NUMERICS Summer School} (2023) ---
        HPC and numerical simulation, CEA
  \item \textbf{IN2P3 School of Statistics} (2022) ---
        Statistical methods for nuclear and particle physics
  \item \textbf{R\,\&\,D Project Management} (2022) --- CEA INSTN
  \item \textbf{Machine Learning and Data Analysis} (2020) ---
        University of Caen
\end{itemize}
}

\cvsection{Technical and Computational Skills}

{\footnotesize
\begin{itemize}[leftmargin=1.4em, itemsep=3pt, topsep=2pt,
                label=\textbullet]
  \item \textbf{Nuclear Modelling:} Langevin stochastic dynamics;
        KEWPIE2/KEWPIE3 framework; cross-section prediction; nuclear
        structure models (collective, mean-field, shell).
  \item \textbf{Uncertainty Quantification and Statistics:}
        Bayesian inference and optimisation; uncertainty
        propagation; model comparison and validation; IN2P3-certified
        statistical methods for nuclear and particle physics.
  \item \textbf{Machine Learning and Data Analysis:} ML for nuclear
        parameter estimation; Python scientific computing; SQL;
        data pipeline development (PostgreSQL, Debezium, Redpanda,
        Spark/Delta Lake, Superset).
  \item \textbf{Scientific Software Development:} Modular
        C\texttt{++}/Python for HPC environments; version-controlled
        workflows (Git); reproducible simulation pipelines.
  \item \textbf{High-Performance Computing:} Coupled differential
        equations; numerical integration and finite-difference schemes;
        HPC workflows (SLURM); sensitivity and convergence analysis.
\end{itemize}
}

\cvsection{Referees}
Available upon request.

\end{document}
